% GENERATED FILE. Edit canonical lessons, then run npm run generate:books.
% canonical-source-hash: fnv1a64:80fd19bfdb50f235
% canonical-lessons: UR-C29-chhe, UR-C29-saat, UR-C29-aath, UR-C29-nau, UR-C29-das, UR-C29-practice

\chapter{Six to Ten, and the Change That Happened Next Door}
\label{ch:ur-count-ten}
\hlchaptermodality{29}

\begin{chapteropening}
\textbf{By the end of this chapter:} \emph{I can count aloud from one to ten in Urdu, and explain why Persian says haft where Urdu says sat.}

Five more numbers with no new letter, the Iranian s-to-h change seen from the side that did not undergo it, and the whole run of ten shuffled.
\end{chapteropening}

\section[chhe]{\ur{چھ} — six, and a puff of air}

\label{lesson:UR-C29-chhe}



\begin{quote}
Count one to five, then five to one.
\end{quote}

\subsection*{You'll want to know first — one word}
\begin{quote}
\textbf{\ur{چھ}} — \emph{chhe} — \textbf{six}
\end{quote}

Two signs and the shortest number in the whole count. The first is the \textbf{\ur{چ}} of \textbf{\ur{چار}}; the second is the two-eyed \textbf{\ur{ھ}} that puts a puff of air on the consonant in front of it.

Hand at the mouth: \emph{chār} moves nothing, \emph{chhe} pushes air out. The tongue does the same thing in both; the lungs do not.

\begin{cousinweb}
\textbf{\ur{چھ}} continues Sanskrit \emph{ṣaṣ}, from Indo-European *\emph{sweks}, \textquotedblleft{}six.\textquotedblright{} English \textbf{six}, Latin \textbf{sex} and Greek \textbf{héx} are the same word.
\end{cousinweb}

\subsection*{Your turn}
\begin{itemize}
\raggedright
  \item \emph{Say it:} \textbf{chhe} — six
  \item \emph{Test:} hand at the mouth — \textbf{chār} still, \textbf{chhe} breathed
  \item \emph{Connect:} \textbf{chhe} $\leftarrow$ Sanskrit \emph{ṣaṣ} $\leftarrow$ *\emph{sweks} $\to$ English \textbf{six}
\end{itemize}

\begin{tcolorbox}[breakable,colback=teal!4,colframe=teal!35!black,title={Before you move on}]
Sources: \href{https://en.wiktionary.org/wiki/\%DA\%86\%DA\%BE}{Wiktionary: \ur{چھ}}; \href{https://www.omniglot.com/language/numbers/urdu.htm}{Omniglot, Numbers in Urdu}.
\end{tcolorbox}
\section[sāt]{\ur{سات} — seven, and a change that happened next door}

\label{lesson:UR-C29-saat}



\begin{quote}
Say \textbf{chhe} and its English cousin.
\end{quote}

\subsection*{You'll want to know first — one word}
\begin{quote}
\textbf{\ur{سات}} — \emph{sāt} — \textbf{seven}
\end{quote}

\begin{cousinweb}
\textbf{\ur{سات}} continues Sanskrit \emph{sapta}, from Indo-European *\emph{septm}, \textquotedblleft{}seven.\textquotedblright{} Latin \textbf{septem}, English \textbf{seven} and Greek \textbf{heptá} are the same word.

Now the neighbour. Persian's seven is \textbf{haft}, and it is the same word — but Iranian turned every Indo-European \emph{s} at the front of a word into an \emph{h}, so \emph{sapta} became \emph{hapta} became \emph{haft}. Indo-Aryan did not, so \emph{sapta} stayed \emph{sāt}.

This is the sharpest illustration in the book of the two roads. A word that came into Urdu FROM Persian carries that change; a word Urdu inherited does not. The numbers are inherited, and the \emph{s} is the proof.
\end{cousinweb}

\subsection*{Your turn}
\begin{itemize}
\raggedright
  \item \emph{Say it:} \textbf{sāt} — seven
  \item \emph{Connect:} \textbf{sāt} $\leftarrow$ Sanskrit \emph{sapta} $\leftarrow$ *\emph{septm} $\to$ Latin \textbf{septem}
  \item \emph{Contrast:} Persian \textbf{haft} — same word, with the \emph{s} turned to \emph{h}
\end{itemize}

\begin{tcolorbox}[breakable,colback=teal!4,colframe=teal!35!black,title={Before you move on}]
Sources: \href{https://en.wiktionary.org/wiki/\%D8\%B3\%D8\%A7\%D8\%AA}{Wiktionary: \ur{سات}}; \href{https://www.omniglot.com/language/numbers/urdu.htm}{Omniglot, Numbers in Urdu}.
\end{tcolorbox}
\section[āṭh]{\ur{آٹھ} — eight, and a sound Europe never had}

\label{lesson:UR-C29-aath}



\begin{quote}
Say \textbf{sāt}, and what Persian did to its first sound.
\end{quote}

\subsection*{You'll want to know first — one word}
\begin{quote}
\textbf{\ur{آٹھ}} — \emph{āṭh} — \textbf{eight}
\end{quote}

Three signs: \textbf{\ur{آ}}, the alif with its wave; \textbf{\ur{ٹ}}, the retroflex; \textbf{\ur{ھ}}, the breath. The middle one is made with the tongue \textbf{curled back} against the roof of the mouth, not against the teeth.

\begin{cousinweb}
\textbf{\ur{آٹھ}} continues Sanskrit \emph{aṣṭa}, from Indo-European *\emph{oktow}, \textquotedblleft{}eight.\textquotedblright{} Latin \textbf{octō}, Greek \textbf{okt\'{\={o}}} and English \textbf{eight} are the same word.

The retroflex is where the family resemblance stops. Latin and English both have a plain \emph{t} or \emph{ct} here; the curled-back consonant is South Asian, and it is one of the sounds that makes an Indo-Aryan language sound like itself rather than like a European one. Eight is a good place to practise it, because the word is short and the sound is stressed.
\end{cousinweb}

\subsection*{Your turn}
\begin{itemize}
\raggedright
  \item \emph{Say it:} \textbf{āṭh} — eight, with the tongue curled back
  \item \emph{Connect:} \textbf{āṭh} $\leftarrow$ Sanskrit \emph{aṣṭa} $\leftarrow$ *\emph{oktow} $\to$ Latin \textbf{octō}
  \item \emph{Run through:} \textbf{chhe, sāt, āṭh}
\end{itemize}

\begin{tcolorbox}[breakable,colback=teal!4,colframe=teal!35!black,title={Before you move on}]
Sources: \href{https://en.wiktionary.org/wiki/\%D8\%A2\%D9\%B9\%DA\%BE}{Wiktionary: \ur{آٹھ}}; \href{https://www.omniglot.com/language/numbers/urdu.htm}{Omniglot, Numbers in Urdu}.
\end{tcolorbox}
\section[nau]{\ur{نو} — nine, and a spelling that carries two words}

\label{lesson:UR-C29-nau}



\begin{quote}
Say \textbf{āṭh}, with the tongue in the right place, and its Latin cousin.
\end{quote}

\subsection*{You'll want to know first — one word}
\begin{quote}
\textbf{\ur{نو}} — \emph{nau} — \textbf{nine}
\end{quote}

Two letters, like the word for two. The second is the same \textbf{\ur{و}} that closes it.

\begin{cousinweb}
\textbf{\ur{نو}} continues Sanskrit \emph{nava}, from Indo-European *\emph{h1newn}, \textquotedblleft{}nine.\textquotedblright{} Latin \textbf{novem}, English \textbf{nine} and Greek \textbf{ennéa} are the same word.

There is a trap here worth meeting once. The same two letters, \textbf{\ur{نو}}, also spell the Persian borrowing \emph{nau}, \textquotedblleft{}new\textquotedblright{} — as in \emph{nau roz}, the new year. That is a DIFFERENT root that happens to arrive at the same spelling by the other road. Only the sentence tells you which is meant.

That is the two roads colliding on one page, and it is the reason this book keeps naming them.
\end{cousinweb}

\subsection*{Your turn}
\begin{itemize}
\raggedright
  \item \emph{Say it:} \textbf{nau} — nine
  \item \emph{Connect:} \textbf{nau} $\leftarrow$ Sanskrit \emph{nava} $\leftarrow$ *\emph{h1newn} $\to$ Latin \textbf{novem}
  \item \emph{Run through:} \textbf{sāt, āṭh, nau}
\end{itemize}

\begin{tcolorbox}[breakable,colback=teal!4,colframe=teal!35!black,title={Before you move on}]
Sources: \href{https://en.wiktionary.org/wiki/\%D9\%86\%D9\%88}{Wiktionary: \ur{نو}}; \href{https://www.omniglot.com/language/numbers/urdu.htm}{Omniglot, Numbers in Urdu}.
\end{tcolorbox}
\section[das]{\ur{دس} — ten}

\label{lesson:UR-C29-das}



\begin{quote}
Say \textbf{nau}, and the other word its two letters can spell.
\end{quote}

\subsection*{You'll want to know first — one word}
\begin{quote}
\textbf{\ur{دس}} — \emph{das} — \textbf{ten}
\end{quote}

Two letters again, and the first is two's \textbf{\ur{د}}.

\begin{cousinweb}
\textbf{\ur{دس}} continues Sanskrit \emph{daśa}, from Indo-European *\emph{dekm}, \textquotedblleft{}ten.\textquotedblright{} Latin \textbf{decem} — and through it \textbf{decimal} and \textbf{December} — English \textbf{ten} and Greek \textbf{déka} are the same word.

\emph{das} is the first number in this book that a shop is likely to say out loud. Very little is priced below ten, which is why a course stopping at five could name amounts nobody charges.
\end{cousinweb}

\subsection*{Your turn}
\begin{itemize}
\raggedright
  \item \emph{Say it:} \textbf{das} — ten
  \item \emph{Connect:} \textbf{das} $\leftarrow$ Sanskrit \emph{daśa} $\leftarrow$ *\emph{dekm} $\to$ Latin \textbf{decem}
  \item \emph{Run through:} \textbf{āṭh, nau, das}
\end{itemize}

\begin{tcolorbox}[breakable,colback=teal!4,colframe=teal!35!black,title={Before you move on}]
Sources: \href{https://en.wiktionary.org/wiki/\%D8\%AF\%D8\%B3}{Wiktionary: \ur{دس}}; \href{https://www.omniglot.com/language/numbers/urdu.htm}{Omniglot, Numbers in Urdu}.
\end{tcolorbox}
\section[Practice]{Payoff — counting aloud to ten}

\label{lesson:UR-C29-practice}



\begin{quote}
Count one to five, then say which of the two roads all five came by.
\end{quote}

\subsection*{Your turn}
\begin{enumerate}
\raggedright
  \item \textbf{Count aloud}, one to ten, without a model — the whole of it, in one breath if you can.
  \item Count backwards from \textbf{das} to \textbf{ek}.
  \item Hear all ten named singly and out of order, and say each.
  \item Read the ten printed shuffled and say each.
  \item Write all ten from dictation. The tranche has added \textbf{no new letter}: every sign in all ten words was already on the page.
\end{enumerate}

Then say the English cousin of each. All ten have one.

\begin{tcolorbox}[breakable,colback=teal!4,colframe=teal!35!black,title={Before you move on}]
Source: \href{https://www.omniglot.com/language/numbers/urdu.htm}{Omniglot, Numbers in Urdu}.
\end{tcolorbox}
